\section{Seven mixed Laurent gates for all three detector orders}

The starting construction is the frozen fifth-round builder, with SHA256
\texttt{68cb3cb1a809766aa102b8c43444d8b38c02760751e773a960489c327e33c12b}.
The selected replacement saves one rational coordinate in each of the three
detectors. It also raises the high detector order from forty-four to
forty-five. The latter is justified by the retained logarithm, not inferred
from a circuit count. The frozen integrated source reviewed here has SHA256
\texttt{3f63cb782d52857446cf39f6e545aeb74f380d1347826eafa888cbb487738d51}.
The inherited global arithmetic qualifications remain unchanged.

\subsection{Use the two independent weight parameters together}

Retain the fifth-round coordinate and guard
\[
 \phi(t)=-49+1/t,\qquad \delta=98t-1,
 \qquad t=A\lambda,
 \qquad \lambda=\frac{32}{1225}+u^2+v^2,
\]
where the three scales are $1,C_0^2,(C_0+C_1)^2$. The preceding
real point bounds give $C_0\ge193$ and $C_0+C_1\ge61$, so every
scale is at least one and is a rational square. The contact polynomial is
\[
 F(t)=t(t-1/49)(t-1/50)(t-1/2450)(t-32/1225).
\]
The required weights and one existing pole-input coordinate are
\[
 e=\frac{(A\delta)^5}{F(t)},\qquad
 q_m=\frac{e^{2m-1}}{\delta^3},\qquad E=e\delta.
\]
In particular $q_1=e/\delta^3$. The order-one product evaluation already
requires $q_1^2,q_1^3,q_1^4$. Move its first two powers, and the existing
coordinate $E$, before the higher-weight circuit. Their computations
are reused by the later product and pole input.

Starting with $e,E,q_1^3$, introduce just seven coordinates:
\begin{align*}
 a&=E^2,& b&=1/a,& c&=b^2/E,& d&=e q_1^3/c,\\
 q_{12}&=d^2/c,& h&=q_{12}^2,& q_{45}&=h^2/q_1^3.
\end{align*}
The exact exponent table is
\[
\begin{array}{c|rrrrrrr}
 &a&b&c&d&q_{12}&h&q_{45}\\ \hline
 e&2&-2&-5&9&23&46&89\\
 \delta&2&-2&-5&-4&-3&-6&-3
\end{array}
\]
For example $c=e^{-5}\delta^{-5}$ and
$d=e\,(e^3\delta^{-9})/(e^{-5}\delta^{-5})=e^9\delta^{-4}$.
Thus the two outputs are exactly the required $q_{12}$ and $q_{45}$.
The old circuit introduced six pure Laurent powers and two output
weights, so this replaces eight new coordinates by seven.

This saving is specific to the actual shared product circuit: counting
$q_1^3$ as free without establishing its existing use would be invalid.
The source checker verifies that the early square and cube are exactly
the aliases later used as the order-one evaluation powers. It also checks
all seven displayed equations against the actual residual rows.
No new use of the auxiliary coordinate $\delta^3$ is introduced; its
three final eliminations remain available.

\subsection{Domains, projections and degrees}

For every solution $t\ge32/1225$ and $\delta\ge39/25>0$.
All factors of $F$ are nonnegative on this interval. The equation
$eF=(A\delta)^5$ rules out $F=0$, so $F,e,E,q_1$ are positive.
Consequently $a,b,c,d,q_{12},h,q_{45}$ are positive and every quotient
denominator in the new chain is nonzero. These coordinates are uniquely
determined by the earlier data. Eliminating them restores the indicated
rational weight identities, while those identities uniquely extend any
point of the preceding domain to the new chain. No roots or branch
choices are introduced.

Each equation has degree at most two in stored rational coordinates:
for instance $cd=e q_1^3$ uses the coordinate already storing $q_1^3$,
and the final relation is $q_1^3q_{45}=h^2$. The existing $E=e\delta$
equation is also quadratic because $\delta$ is affine in $t$.
Reordering it does not create a dependency cycle. The source checker
checks the degree of every source residual with the fixed anchor
constant represented by a coordinate.

The total-positive factors $1+q_mD_j^2$ remain at least one at both real
embeddings. Their product $W=W_0+W_1\sqrt2$ therefore has $W_0\ge1$
and $N(W)>0$. The actual pole input still is
\[
 B=\frac{\bigl(N(W)/(2W_0)\bigr)^2}{E}.
\]
Thus the trace quotient and pole denominator have the same nonzero
domains as before. There is no new sign predicate.

\subsection{The arbitrary-order local consumer}

The sixteen field tests now have orders $1,12,45$: nine order-one tests,
one input cross test of order twelve, and six high-order tests. The latter
are the primary inverse root, the three Lucas-label cross tests, the input
discrepancy, and the direct Lucas residual. The product grouping remains
one mixed $(1,12)$ block, two four-test order-one blocks, and two
three-test high-order blocks. The mixed-block identity is unchanged:
\[
 (1+q_1U)(1+q_{12}V)
 =1+q_1U+q_{12}V+(q_1q_{12})UV.
\]

Put $R=E^{-1}$. For every positive integer $m$,
\[
 1+R\left(\frac{D_j/\delta^{m+1}}{R^m}\right)^2
   =1+q_mD_j^2.
\]
The source's arbitrary-order one-place argument therefore applies
without changing its valuation calculation. More explicitly, at a
consumed unramified prime with odd $\alpha=v_p(R)$, the nonconstant term
of each factor has odd valuation. Thus the factor valuation is its
minimum with zero. The pole condition gives positive valuation to
$R/\operatorname{Tr}(W^{-1})^2$, forcing $\alpha>0$ and, at a common
successful place, $v(D_j)>(m-3/4)\alpha$. Integral normalized valuations
give $v(D_j)\ge m$. The guard is a unit at the consumed contacts, as
already established for the retained exceptional set. This is the actual
proof of Lemma \texttt{v5-one-place} in the pinned upstream manuscript,
lines 11031--11061; its argument permits arbitrary positive orders.
The reduced exceptional set and pole hypotheses are justified in the
previous arithmetic proofs, rather than by disregarding the printed
large prime cutoff.

The retained long logarithm polynomial is $G_{12}$. If a formal
parameter has depth $r\ge1$, its polynomial argument has valuation $4r$.
For every omitted coefficient index $j\ge13$ and every consumed prime
$p\ge11$, the existing coefficient estimate gives
\[
 v_p(a_j)\ge-\frac{4j+1}{p-1},\qquad
 \frac{4j+1}{p-1}\le\frac{4j+1}{10}\le4j-45.
\]
The last inequality follows from
$10(4j-45)-(4j+1)=36(j-13)+17\ge0$.
Hence every omitted term has valuation at least $45r$, and the valuations
tend to infinity. This proves the required relative precision forty-five
without adding an exceptional prime. In particular no unproved uniform
precision fifty-two is used. The input logarithm keeps $G_2$ and precision
twelve. The primary inverse-root test gives primary depth at least
forty-five, and hence primary relative error at least one hundred eighty.
The previously proved weak-unit comparison first transfers integrality;
then the two-anchor comparison gives the actual Lucas-label congruences
at order forty-five and the input-label congruence at order twelve.
The separate identity-input branch is retained.

Completeness replaces only the primary index progression $\Pi^{43}\mid h$
by $\Pi^{44}\mid h$. The original separation argument gives
$\gcd(\Pi,M_+L)=1$, so this change is compatible with the same CRT
requirements. The alternate points still require only formal depth one;
there is no incompatible demand that both $h$ and $h-1$ be divisible by
the chosen prime. All finite lower bounds on prepared point functions
remain available.

The fifth-round real and small-prime limits are uniform in every fixed
positive order. In particular the nonconstant factors are
$O(\lambda^{-3})$ at infinity. At $p=2,3,17$, writing
$\nu=v_p(v)$ for the chosen denominator parameter,
$\eta=v_p(A)$ and $\alpha=v_p(98)$, their lower bound is
\[
 3\nu+(10m-8)(\eta+\alpha)-2c_p.
\]
It tends to infinity for $m=1,12,45$. Only finite preparatory constants
change. The fixed Gaussian squareclasses and all contact values are
unchanged. This verifies the real and local completeness transfer;
the seven extra rational coordinates themselves impose no approximation
or smallness condition.

\subsection{Implementation and exact measurements}

The isolated detector replacement was compiled completely:
\[
 412/380\quad\longrightarrow\quad402/370
 \quad\longrightarrow\quad398/366.
\]
These are coordinates/residuals before compilation, after the generic
compiler, and after four further certified eliminations. The baseline
was $415/383\to405/373\to401/369$. Thus all three raw-coordinate savings
survive the final projection-preserving elimination pipeline.

Combined with the independently proved two-coordinate scalar norm-tree
reduction and the inverse-root physical anchor, the selected source has
\[
 410/378\quad\longrightarrow\quad400/368
 \quad\longrightarrow\quad\boxed{396/364}.
\]
The inverse-root anchor uses $S_0=3^{3072}$; the exact anchor equivalence
and scalar/Lucas height margins are proved in
\texttt{arithmetic\_round6.tex}. The mixed scalar norm-tree proof is
\texttt{geometry\_round6.tex}. Neither change is inferred merely from
matching counts. The final quartic has 396 active rational coordinates,
364 quadratic residuals, 26,225 ordinary-coordinate monomials and
primitive integer content one. The three ordinary decimal constant
circuits have respectively 398, 399 and 400 active coordinates. Their
coefficient and support claims belong to the separate exact table audit.

The selected checker \texttt{check\_detector\_round6.py} imports neither
producer nor compiler. It passes 390 source assertions, seven exact
Laurent exponent identities and two expanded mixed-block identities on
the frozen integrated source. Its coverage includes independent binomial
reconstruction of the $G_{12}$ square completion, all logarithm outputs
and quotient consumers, every field-test alias, all squared tests and
shared symmetric coefficients, all evaluation powers, the shortened
scales and guard, both high-weight outputs, actual square/cube reuse,
the norm and pole input, and the lifted degree bound. The independent
scalar reviewer additionally checks 64 actual source equations, seven
chart/proxy identities and the exact valuation descent. These are finite
algebraic checks; the global arithmetic inputs retain their stated
conditional status.

\subsection{Search limitations and unselected candidates}

The circuit was found by an exact exponent-pair search using operations
$(a,b)+(c,d)-(e,f)$, corresponding to quadratic quotient gates. The
first three new operations were restricted to doubled differences;
the final two used general pair sums minus one earlier exponent pair.
The search included the already required $e,\delta,\delta^2,\delta^3,E$
and order-one powers. The selected search additionally prohibited new
uses of $\delta^3$ to preserve its existing eliminations. This is a
restricted constructive search, not an exhaustive lower-bound proof
for every rational circuit. No minimality assertion is made.

Cost-neutral order-forty-six and order-forty-eight circuits were also
fully compiled; both retained 401 final coordinates before the separate
scalar improvement. Seven-gate alternatives using $\delta^3$ and lower
orders ten or eleven were investigated, but were not selected over the
simple measured order-twelve/forty-five circuit. No unselected variant
is used by the final arithmetic argument.
